\documentclass[14pt]{extarticle}
% =============================================================================
% Shared Preamble for Worksheets, Homeworks, Exams & Lesson Notes
% (scripts/tex_preamble.tex)
% =============================================================================

\newif\ifdetailed
\ifdefined\setdetailed
  \setdetailed
\fi

\newif\ifwithsols
\ifdefined\setwithsols
  \setwithsols
\fi

\newif\ifwhiteblack
\ifdefined\setwhiteblack
  \setwhiteblack
\fi

% Document-class specific packages (loaded before common_preamble.tex / polyglossia)
\usepackage{fancyhdr}
\usepackage{enumitem}
\usepackage{titlesec}
\usepackage[bottom]{footmisc}
\usepackage{titling}
\usepackage{longtable} % page-breakable tables
\usepackage{float}
\usepackage[onehalfspacing]{setspace}

\ifwhiteblack
  \usepackage[a4paper, landscape, top=1.25cm, bottom=1.25cm, left=1.25cm, right=1.25cm, headheight=1.5cm, headsep=0.5cm]{geometry}
\else
  \usepackage[a4paper, top=2.4cm, bottom=2.4cm, left=2.2cm, right=2.2cm, headheight=1.4cm, headsep=0.4cm, footskip=1.0cm]{geometry}
\fi

% Shared common packages, fonts, polyglossia, callout boxes, code helpers, date & header macros
\IfFileExists{scripts/common_preamble.tex}{%
  % =============================================================================
% Shared Common Preamble (scripts/common_preamble.tex)
% Included by:
%   - scripts/tex_preamble.tex    (worksheets, homeworks, exams, lesson notes)
%   - scripts/book_preamble.tex   (two-sided Arabic book)
%   - scripts/beamer_preamble.tex (Beamer presentation slides)
% =============================================================================

\usepackage{fontspec}
\usepackage{amsmath}
\usepackage{amssymb}
\usepackage{xcolor}
\usepackage{graphicx}
\usepackage{listings}
\usepackage[most]{tcolorbox}
\usepackage{pifont}
\usepackage{minted}
\usepackage{etoolbox}
\usepackage{array}
\usepackage{booktabs}
\usepackage{tikz}
\usepackage{datetime2}

\usetikzlibrary{automata, positioning, arrows.meta, arrows, shapes.geometric, calc}

% Beamer already loads hyperref internally; load it here for article/book classes
% right before polyglossia (bidi).
\makeatletter
\@ifclassloaded{beamer}{}{%
  \usepackage{hyperref}%
  \hypersetup{hidelinks}%
  \providecommand{\LTcaptype}{table}%
}
\makeatother

% Load polyglossia (and bidi) AFTER graphics, tables, boxes, code, and hyperref
\usepackage{polyglossia}
\setmainlanguage[numerals=western]{arabic}
\setotherlanguage{english}
\newfontfamily\arabicfont[Script=Arabic]{Calibri}
\newfontfamily\arabicfontsf[Script=Arabic]{Calibri}
\newfontfamily\arabicfonttt[Script=Arabic]{Courier New}
\newfontfamily\emoji[BoldFont=*,ItalicFont=*]{Segoe UI Emoji}

% Allow slight inter-word stretch to avoid Overfull \hbox with inline \textenglish
\emergencystretch=1.5em

% -----------------------------------------------------------------------------
% Default Minted & Callout Box Definitions
% -----------------------------------------------------------------------------
\AtBeginEnvironment{minted}{\setstretch{0.95}}
\setminted{fontsize=\small, breaklines=true, vspace=2pt}

\tcbset{
  arc=2.5pt, outer arc=2.5pt,
  boxsep=1pt, left=6pt, right=6pt, top=3pt, bottom=3pt,
  toptitle=1.5pt, bottomtitle=1.5pt,
  before skip=6pt, after skip=6pt,
  fonttitle=\bfseries,
  fontupper=\setstretch{1.1}\setlength{\parskip}{2pt},
  colback=white, colframe=black, boxrule=0.6pt,
  boxcolor/.style 2 args={colback=#1, colframe=#2}
}

\newtcolorbox{boxDef}[1][]{boxcolor={blue!4!white}{blue!75!black},
  title={{\emoji📘} تعريف\ifx\\#1\\\else ~#1\fi :}}
\newtcolorbox{boxExercise}[1][]{boxcolor={cyan!4!white}{cyan!70!black},
  title={{\emoji🧩} تمرين\ifx\\#1\\\else ~#1\fi :}}
\newtcolorbox{boxExample}[1][]{boxcolor={yellow!6!white}{orange!85!black},
  title={{\emoji📝} مثال\ifx\\#1\\\else ~#1\fi :}}
\newtcolorbox{boxNote}[1][]{boxcolor={gray!7!white}{black!75},
  title={{\emoji✨} ملاحظة\ifx\\#1\\\else ~#1\fi :}}
\newtcolorbox{boxAttention}[1][]{boxcolor={magenta!5!white}{magenta!75!black},
  title={{\emoji🔔} تنبيه\ifx\\#1\\\else ~#1\fi :}}
\newtcolorbox{boxWarning}[1][]{boxcolor={red!4!white}{red!75!black},
  title={{\emoji⚡} ملاحظة هامة\ifx\\#1\\\else ~#1\fi :}}
\newtcolorbox{boxSolution}[1][]{boxcolor={green!4!white}{green!60!black},
  title={{\emoji✅} حل\ifx\\#1\\\else ~#1\fi :}}
\newtcolorbox{boxSymbol}[1][]{boxcolor={purple!4!white}{purple!70!black},
  title={{\emoji🔣} رمز\ifx\\#1\\\else ~#1\fi :}}
\newtcolorbox{boxHint}[1][]{boxcolor={teal!4!white}{teal!65!black},
  title={{\emoji💡} تلميح\ifx\\#1\\\else ~#1\fi :}}

% Compact code boxes: tighter line spacing and small vertical gaps
\tcbset{
  simplecode/.style={
    arc=2pt, outer arc=2pt,
    left=4pt, right=4pt, top=4pt, bottom=4pt,
    fontupper=\setstretch{0.95},
    before skip=6pt, after skip=6pt,
    colback=gray!5, colframe=black!50, boxrule=0.4pt
  }
}
\newenvironment{boxCode}{\begin{tcolorbox}[simplecode]}{\end{tcolorbox}}

% Code cell used inside tables (keeps code left-to-right inside RTL tables)
% Use inside a fixed-width p{...} column: \begin{pycell}\begin{minted}...\end{minted}\end{pycell}
\makeatletter
\newtcolorbox{pycell}{frame empty,colback=white,boxsep=0pt,left=2pt,right=2pt,top=2pt,bottom=2pt,width=\linewidth,before=\noindent,after=,baseline={\tcb@height-\ht\strutbox},fontupper=\setstretch{0.95}\setlength{\topsep}{0pt}\setlength{\partopsep}{0pt}\setlength{\parskip}{0pt}}
\makeatother

\newcolumntype{C}[1]{>{\centering\arraybackslash}p{#1}}

% Dotted date style (DD.MM.YYYY)
\makeatletter
\DTMnewdatestyle{dotted}{%
  \renewcommand{\DTMdisplaydate}[4]{\two@digits{##3}.\two@digits{##2}.##1}%
  \renewcommand{\DTMDisplaydate}{\DTMdisplaydate}%
}
\makeatother
\DTMsetdatestyle{dotted}

% -----------------------------------------------------------------------------
% Customizable Header / Footer Metadata
% -----------------------------------------------------------------------------
\newcommand{\headerleft}{مدرسة التسامح الشاملة}
\newcommand{\headerright}{الأستاذ محمود اغبارية}
\newcommand{\footercenter}{\thepage}

\newcommand{\setheaderleft}[1]{\renewcommand{\headerleft}{#1}}
\newcommand{\setheaderright}[1]{\renewcommand{\headerright}{#1}}
\newcommand{\setfootercenter}[1]{\renewcommand{\footercenter}{#1}}
%
}{%
  \IfFileExists{../../../scripts/common_preamble.tex}{%
    % =============================================================================
% Shared Common Preamble (scripts/common_preamble.tex)
% Included by:
%   - scripts/tex_preamble.tex    (worksheets, homeworks, exams, lesson notes)
%   - scripts/book_preamble.tex   (two-sided Arabic book)
%   - scripts/beamer_preamble.tex (Beamer presentation slides)
% =============================================================================

\usepackage{fontspec}
\usepackage{amsmath}
\usepackage{amssymb}
\usepackage{xcolor}
\usepackage{graphicx}
\usepackage{listings}
\usepackage[most]{tcolorbox}
\usepackage{pifont}
\usepackage{minted}
\usepackage{etoolbox}
\usepackage{array}
\usepackage{booktabs}
\usepackage{tikz}
\usepackage{datetime2}

\usetikzlibrary{automata, positioning, arrows.meta, arrows, shapes.geometric, calc}

% Beamer already loads hyperref internally; load it here for article/book classes
% right before polyglossia (bidi).
\makeatletter
\@ifclassloaded{beamer}{}{%
  \usepackage{hyperref}%
  \hypersetup{hidelinks}%
  \providecommand{\LTcaptype}{table}%
}
\makeatother

% Load polyglossia (and bidi) AFTER graphics, tables, boxes, code, and hyperref
\usepackage{polyglossia}
\setmainlanguage[numerals=western]{arabic}
\setotherlanguage{english}
\newfontfamily\arabicfont[Script=Arabic]{Calibri}
\newfontfamily\arabicfontsf[Script=Arabic]{Calibri}
\newfontfamily\arabicfonttt[Script=Arabic]{Courier New}
\newfontfamily\emoji[BoldFont=*,ItalicFont=*]{Segoe UI Emoji}

% Allow slight inter-word stretch to avoid Overfull \hbox with inline \textenglish
\emergencystretch=1.5em

% -----------------------------------------------------------------------------
% Default Minted & Callout Box Definitions
% -----------------------------------------------------------------------------
\AtBeginEnvironment{minted}{\setstretch{0.95}}
\setminted{fontsize=\small, breaklines=true, vspace=2pt}

\tcbset{
  arc=2.5pt, outer arc=2.5pt,
  boxsep=1pt, left=6pt, right=6pt, top=3pt, bottom=3pt,
  toptitle=1.5pt, bottomtitle=1.5pt,
  before skip=6pt, after skip=6pt,
  fonttitle=\bfseries,
  fontupper=\setstretch{1.1}\setlength{\parskip}{2pt},
  colback=white, colframe=black, boxrule=0.6pt,
  boxcolor/.style 2 args={colback=#1, colframe=#2}
}

\newtcolorbox{boxDef}[1][]{boxcolor={blue!4!white}{blue!75!black},
  title={{\emoji📘} تعريف\ifx\\#1\\\else ~#1\fi :}}
\newtcolorbox{boxExercise}[1][]{boxcolor={cyan!4!white}{cyan!70!black},
  title={{\emoji🧩} تمرين\ifx\\#1\\\else ~#1\fi :}}
\newtcolorbox{boxExample}[1][]{boxcolor={yellow!6!white}{orange!85!black},
  title={{\emoji📝} مثال\ifx\\#1\\\else ~#1\fi :}}
\newtcolorbox{boxNote}[1][]{boxcolor={gray!7!white}{black!75},
  title={{\emoji✨} ملاحظة\ifx\\#1\\\else ~#1\fi :}}
\newtcolorbox{boxAttention}[1][]{boxcolor={magenta!5!white}{magenta!75!black},
  title={{\emoji🔔} تنبيه\ifx\\#1\\\else ~#1\fi :}}
\newtcolorbox{boxWarning}[1][]{boxcolor={red!4!white}{red!75!black},
  title={{\emoji⚡} ملاحظة هامة\ifx\\#1\\\else ~#1\fi :}}
\newtcolorbox{boxSolution}[1][]{boxcolor={green!4!white}{green!60!black},
  title={{\emoji✅} حل\ifx\\#1\\\else ~#1\fi :}}
\newtcolorbox{boxSymbol}[1][]{boxcolor={purple!4!white}{purple!70!black},
  title={{\emoji🔣} رمز\ifx\\#1\\\else ~#1\fi :}}
\newtcolorbox{boxHint}[1][]{boxcolor={teal!4!white}{teal!65!black},
  title={{\emoji💡} تلميح\ifx\\#1\\\else ~#1\fi :}}

% Compact code boxes: tighter line spacing and small vertical gaps
\tcbset{
  simplecode/.style={
    arc=2pt, outer arc=2pt,
    left=4pt, right=4pt, top=4pt, bottom=4pt,
    fontupper=\setstretch{0.95},
    before skip=6pt, after skip=6pt,
    colback=gray!5, colframe=black!50, boxrule=0.4pt
  }
}
\newenvironment{boxCode}{\begin{tcolorbox}[simplecode]}{\end{tcolorbox}}

% Code cell used inside tables (keeps code left-to-right inside RTL tables)
% Use inside a fixed-width p{...} column: \begin{pycell}\begin{minted}...\end{minted}\end{pycell}
\makeatletter
\newtcolorbox{pycell}{frame empty,colback=white,boxsep=0pt,left=2pt,right=2pt,top=2pt,bottom=2pt,width=\linewidth,before=\noindent,after=,baseline={\tcb@height-\ht\strutbox},fontupper=\setstretch{0.95}\setlength{\topsep}{0pt}\setlength{\partopsep}{0pt}\setlength{\parskip}{0pt}}
\makeatother

\newcolumntype{C}[1]{>{\centering\arraybackslash}p{#1}}

% Dotted date style (DD.MM.YYYY)
\makeatletter
\DTMnewdatestyle{dotted}{%
  \renewcommand{\DTMdisplaydate}[4]{\two@digits{##3}.\two@digits{##2}.##1}%
  \renewcommand{\DTMDisplaydate}{\DTMdisplaydate}%
}
\makeatother
\DTMsetdatestyle{dotted}

% -----------------------------------------------------------------------------
% Customizable Header / Footer Metadata
% -----------------------------------------------------------------------------
\newcommand{\headerleft}{مدرسة التسامح الشاملة}
\newcommand{\headerright}{الأستاذ محمود اغبارية}
\newcommand{\footercenter}{\thepage}

\newcommand{\setheaderleft}[1]{\renewcommand{\headerleft}{#1}}
\newcommand{\setheaderright}[1]{\renewcommand{\headerright}{#1}}
\newcommand{\setfootercenter}[1]{\renewcommand{\footercenter}{#1}}
%
  }{%
    \input{../../scripts/common_preamble.tex}%
  }%
}

\pagestyle{fancy}

% Prevent widow/orphan lines
\widowpenalty=10000
\clubpenalty=10000

% Compact section headings and default list spacing so pages do not waste vertical space
\titlespacing*{\section}{0pt}{1.4ex plus 0.4ex minus 0.2ex}{0.7ex plus 0.2ex}
\titlespacing*{\subsection}{0pt}{1.1ex plus 0.3ex minus 0.2ex}{0.5ex plus 0.1ex}
\titlespacing*{\subsubsection}{0pt}{0.9ex plus 0.2ex minus 0.1ex}{0.4ex plus 0.1ex}
\setlist{topsep=4pt, partopsep=0pt, parsep=2pt, itemsep=4pt}

\gappto\captionsarabic{\renewcommand{\contentsname}{جدول المحتويات}}

% Legacy listings settings
\lstset{
  language=[Sharp]C,
  numbers=left,
  stepnumber=1,
  numbersep=8pt,
  frame=single,
  basicstyle=\ttfamily\small
}
\ifwhiteblack
  \lstset{
    keywordstyle=\bfseries,
    stringstyle=\ttfamily,
    commentstyle=\itshape
  }
\else
  \lstset{
    keywordstyle=\color{blue},
    stringstyle=\color{red},
    commentstyle=\color{green!50!black}
  }
\fi

% White-and-black landscape presentation mode overrides
\ifwhiteblack
  \AtBeginDocument{
    \LARGE
    \usemintedstyle{bw}
  }
  \setminted{vspace=2pt}
  \tcbset{
    colback=white, colframe=black, boxrule=1pt,
    boxcolor/.style 2 args={colback=white, colframe=black},
    simplecode/.append style={colback=white, colframe=black, boxrule=1pt}
  }
\fi

% Default to empty author to avoid "LaTeX Warning: No \author given" when omitted
\author{}

% Compact title block
\makeatletter
\renewcommand{\@maketitle}{%
  \begin{center}%
    {\ifwhiteblack\Huge\else\LARGE\fi \bfseries \@title \par}%
    \ifdefempty{\@author}{}{%
      \vskip 0.15em%
      {\large \@author \par}%
    }%
    {\small \hfill \textenglish{\DTMtoday} \par}%
  \end{center}%
  \vspace{-0.3cm}%
}
\makeatother

\apptocmd{\maketitle}{\thispagestyle{fancy}}{}{}

% Apply page headers and footer (using \headerleft, \headerright, \footercenter from common_preamble.tex)
\newcommand{\applyfancyheaders}{%
  \fancyhf{}%
  \fancyhead[L]{\textarabic{\headerleft}}%
  \fancyhead[R]{\textarabic{\headerright}}%
  \fancyfoot[C]{\textarabic{\footercenter}}%
}

\fancypagestyle{plain}{\applyfancyheaders}
\applyfancyheaders

\setcounter{tocdepth}{3} % only section, subsection, and subsubsection in TOC

\newcommand{\importpdfpage}[2]{%
    \newgeometry{left=0cm,right=0cm,top=0cm,bottom=0cm}
    \thispagestyle{empty}
    \includegraphics[page=#2,width=0.95\paperwidth,height=0.95\paperheight,keepaspectratio]{#1}
    \restoregeometry
}

\newcommand{\insertFullImg}[1]{%
    \noindent
    \makebox[\textwidth][c]{
        \includegraphics[width=0.9\paperwidth,keepaspectratio]{#1}%
    }%
}%


\ifwithsols
\title{حل ورقة تمرن $4$ للصف العاشر $10$ - المتغيرات ومكتبة \textenglish{Math}}
\else
\title{ورقة تمرن $4$ للصف العاشر $10$ - المتغيرات ومكتبة \textenglish{Math}}
\fi

\begin{document}

\maketitle
\thispagestyle{fancy}

\begin{enumerate}[itemsep=2em]

% ==============================================================
% السؤال 1: تتبّع بدون جملة شرط (متغيرات وحسابات وعمليات Math)
% ==============================================================
\item
أمامك مقطع البرنامج التالي الذي يستقبل عددًا صحيحًا موجبًا \textenglish{num} مكوّنًا من $3$ منازل (بين $100$ و $999$):

\begin{boxCode}
\begin{english}
\begin{minted}{csharp}
int num = int.Parse(Console.ReadLine());
int a = num / 100;
int b = (num / 10) % 10;
int c = num % 10;
int m = Math.Max(a, c);
int d = Math.Abs(a - b);
int res = m * 10 + d;
Console.WriteLine(res);
\end{minted}
\end{english}
\end{boxCode}

\begin{enumerate}[itemsep=1.5em]
    \filbreak
    \item \textbf{ما هو المُخرج؟} \\
    تتبّع تنفيذ المقطع البرمجي أعلاه لكل قيمة من قيم المُدخل \textenglish{num} في الجدول التالي، واكتب قيم المتغيرات والمُخرج الذي يطبعه البرنامج:
    \begin{center}
    \begin{english}
    \renewcommand{\arraystretch}{1.45}
    \begin{tabular}{|C{1.8cm}|C{1.2cm}|C{1.2cm}|C{1.2cm}|C{1.2cm}|C{1.2cm}|C{2.2cm}|}
    \hline
    \textbf{num} & \textbf{a} & \textbf{b} & \textbf{c} & \textbf{m} & \textbf{d} & \textarabic{\textbf{المُخرج}} \\
    \hline
    472 & & & & & & \\
    \hline
    389 & & & & & & \\
    \hline
    661 & & & & & & \\
    \hline
    515 & & & & & & \\
    \hline
    \end{tabular}
    \end{english}
    \end{center}
    \filbreak

    \item \textbf{أعطِ مُدخلاً يطبع كلّ واحد من هذه المخرجات:} \\
    لكل مُخرج في الجدول التالي، اقترح عددًا صحيحًا \textenglish{num} مكوّنًا من $3$ منازل (\textbf{مختلفًا} عن أعداد الفرع أ) يجعل البرنامج يطبع هذا المُخرج، مع تحقيق الشرط الإضافي المذكور:
    \begin{center}
    \renewcommand{\arraystretch}{1.45}
    \begin{tabular}{|C{2.4cm}|p{7cm}|C{2.4cm}|}
    \hline
    \textbf{المُخرج المطلوب} & \textbf{شرط إضافي على المُدخلات} & \textbf{قيمة \textenglish{num}} \\
    \hline
    $83$ & بحيث تكون منزلة المئات أكبر من منزلة الآحاد ($a > c$) & \\
    \hline
    $83$ & بحيث تكون منزلة المئات أصغر من منزلة الآحاد ($a < c$) & \\
    \hline
    $60$ & بحيث تكون جميع منازل العدد غير صفرية & \\
    \hline
    $75$ & بحيث تكون منزلة المئات أصغر من منزلة العشرات ($a < b$) & \\
    \hline
    \end{tabular}
    \end{center}
    \filbreak
\end{enumerate}

\ifwithsols
\begin{boxSolution}[الفرع أ والفرع ب]
\textbf{حل الفرع أ:}
\begin{center}
\begin{english}
\renewcommand{\arraystretch}{1.35}
\begin{tabular}{|C{1.8cm}|C{1.2cm}|C{1.2cm}|C{1.2cm}|C{1.2cm}|C{1.2cm}|C{2.2cm}|}
\hline
\textbf{num} & \textbf{a} & \textbf{b} & \textbf{c} & \textbf{m} & \textbf{d} & \textarabic{\textbf{المُخرج}} \\
\hline
472 & 4 & 7 & 2 & 4 & 3 & 43 \\
\hline
389 & 3 & 8 & 9 & 9 & 5 & 95 \\
\hline
661 & 6 & 6 & 1 & 6 & 0 & 60 \\
\hline
515 & 5 & 1 & 5 & 5 & 4 & 54 \\
\hline
\end{tabular}
\end{english}
\end{center}

\textbf{حل الفرع ب (أمثلة لإجابات صحيحة):}
\begin{itemize}[nosep]
    \item لطباعة $83$ مع $a > c$: يجب أن يكون $m = 8$ و $d = |a - b| = 3$، وبما أن $a > c$ فإن $a = 8$ و $b = 5$ و $c < 8$، مثل: $\text{num} = 852$.
    \item لطباعة $83$ مع $a < c$: يجب أن يكون $c = 8$ و $a < 8$ و $|a - b| = 3$، مثل: $\text{num} = 478$ (أو $528$).
    \item لطباعة $60$ بمنازل غير صفرية: يجب أن يكون $m = 6$ و $d = 0$ (أي $a = b$)، مثل: $\text{num} = 664$ (أو $336$).
    \item لطباعة $75$ مع $a < b$: يجب أن يكون $m = 7$ و $b - a = 5$، مثل: $\text{num} = 277$ (أو $387$ أو $497$).
\end{itemize}
\end{boxSolution}
\fi
\filbreak

% ==============================================================
% الأسئلة 2-4: حساب تعبيرات رياضية باستخدام مكتبة Math
% ==============================================================
\item
اكتب برنامجًا يستقبل من المستخدم عددًا عشريًّا يمثّل قيمة $x$، ويحسب ويطبع قيمة التعبير التالي عند النقطة $x$ (افترض أنّ المدخلات دائمًا صحيحة ولا تسبب خطأ رياضيًّا):
\[ f(x) = \frac{x + 1}{x^3} \]
\ifdetailed
\begin{boxExample}
\begin{english}
\begin{minted}{text}
Please enter x:
2
f(x) = 0.375
\end{minted}
\end{english}
\end{boxExample}
\fi
\ifwithsols
\begin{boxSolution}
\begin{english}
\begin{minted}{csharp}
double x = double.Parse(Console.ReadLine());

double num = x + 1;
double den = Math.Pow(x, 3);
double result = num / den;

Console.WriteLine("f(x) = " + result);
\end{minted}
\end{english}
\end{boxSolution}
\fi
\filbreak

\item
اكتب برنامجًا يستقبل من المستخدم عددًا عشريًّا يمثّل قيمة $x$، ويحسب ويطبع قيمة التعبير التالي عند النقطة $x$ (افترض أنّ المدخلات دائمًا صحيحة ولا تسبب خطأ رياضيًّا):
\[ f(x) = (\sqrt{x} + 1)^2 \]
\ifdetailed
\begin{boxExample}
\begin{english}
\begin{minted}{text}
Please enter x:
4
f(x) = 9
\end{minted}
\end{english}
\end{boxExample}
\fi
\ifwithsols
\begin{boxSolution}
\begin{english}
\begin{minted}{csharp}
double x = double.Parse(Console.ReadLine());

double root = Math.Sqrt(x);
double inner = root + 1;
double result = Math.Pow(inner, 2);

Console.WriteLine("f(x) = " + result);
\end{minted}
\end{english}
\end{boxSolution}
\fi
\filbreak

\item
اكتب برنامجًا يستقبل من المستخدم عددًا عشريًّا يمثّل قيمة $x$، ويحسب ويطبع قيمة التعبير التالي عند النقطة $x$ (افترض أنّ المدخلات دائمًا صحيحة ولا تسبب خطأ رياضيًّا):
\[ f(x) = \frac{|x^3 - 8|}{\sqrt{x^2 + 16}} \]
\ifdetailed
\begin{boxExample}
\begin{english}
\begin{minted}{text}
Please enter x:
3
f(x) = 3.8
\end{minted}
\end{english}
\end{boxExample}
\fi
\ifwithsols
\begin{boxSolution}
\begin{english}
\begin{minted}{csharp}
double x = double.Parse(Console.ReadLine());

double num = Math.Abs(Math.Pow(x, 3) - 8);
double den = Math.Sqrt(Math.Pow(x, 2) + 16);
double result = num / den;

Console.WriteLine("f(x) = " + result);
\end{minted}
\end{english}
\end{boxSolution}
\fi
\filbreak

% ==============================================================
% السؤال 5: مساحة المربع ونصف قطر الدائرة مع Math.Round
% ==============================================================
\item
اكتب برنامجًا يستقبل محيط مربع، ويحسب ويطبع مساحة هذا المربع. \\
ثم، بافتراض وجود دائرة لها \textbf{نفس مساحة هذا المربع}، يحسب البرنامج نصف قطر هذه الدائرة ويقرّبه لمنزلتين عشريتين باستخدام العملية \textenglish{Math.Round}، ويطبعه. \\
($\text{المساحة} = 3.14 \times r^2 \Rightarrow r = \sqrt{\frac{\text{المساحة}}{3.14}}$)
\ifdetailed
\begin{boxExample}
\begin{english}
\begin{minted}{text}
Please enter the square's perimeter:
20
Square Area = 25
Circle Radius (rounded) = 2.82
\end{minted}
\end{english}
\end{boxExample}
\fi
\ifwithsols
\begin{boxSolution}
\begin{english}
\begin{minted}{csharp}
double perimeter = double.Parse(Console.ReadLine());

double side = perimeter / 4.0;
double squareArea = Math.Pow(side, 2);
Console.WriteLine("Square Area = " + squareArea);

double radiusSq = squareArea / 3.14;
double radius = Math.Sqrt(radiusSq);
radius = Math.Round(radius, 2);

Console.WriteLine("Circle Radius (rounded) = " + radius);
\end{minted}
\end{english}
\end{boxSolution}
\fi
\filbreak

% ==============================================================
% السؤال 6: تفكيك المبلغ إلى أوراق نقدية
% ==============================================================
\item
اكتب برنامجًا يستقبل من المستخدم عددًا صحيحًا يمثّل مبلغًا من المال (من مضاعفات الـ $50$)، ويطبع أقل عدد ممكن من الأوراق النقدية (من فئات $200$، $100$، $50$) اللازمة لتكوين هذا المبلغ.
\ifdetailed
\begin{boxExample}
\begin{english}
\begin{minted}{text}
Please enter the amount:
850
200 bills: 4
100 bills: 0
50 bills: 1
\end{minted}
\end{english}
\end{boxExample}
\fi
\ifwithsols
\begin{boxSolution}
\begin{english}
\begin{minted}{csharp}
int amount = int.Parse(Console.ReadLine());

int bills200 = amount / 200;
amount = amount % 200;

int bills100 = amount / 100;
amount = amount % 100;

int bills50 = amount / 50;

Console.WriteLine("200 bills: " + bills200);
Console.WriteLine("100 bills: " + bills100);
Console.WriteLine("50 bills: " + bills50);
\end{minted}
\end{english}
\end{boxSolution}
\fi
\filbreak

% ==============================================================
% السؤال 7: عكس عدد من 3 منازل والفارق المطلق
% ==============================================================
\item
اكتب برنامجًا يستقبل من المستخدم عددًا صحيحًا من $3$ منازل (مثلاً $258$)، ويقوم ببناء عدد جديد بحيث تكون منازله معكوسة ($852$). ثم يطبع العدد الجديد، ويحسب ويطبع الفارق المطلق بين العدد الأصلي والعدد المعكوس.
\ifdetailed
\begin{boxExample}
\begin{english}
\begin{minted}{text}
Please enter a 3-digit number:
258
Reversed number = 852
Absolute difference = 594
\end{minted}
\end{english}
\end{boxExample}
\fi
\ifwithsols
\begin{boxSolution}
\begin{english}
\begin{minted}{csharp}
int num = int.Parse(Console.ReadLine());
int units = num % 10;
int tens = (num / 10) % 10;
int hundreds = num / 100;

int reversedNum = units * 100 + tens * 10 + hundreds;
int diff = Math.Abs(num - reversedNum);

Console.WriteLine("Reversed number = " + reversedNum);
Console.WriteLine("Absolute difference = " + diff);
\end{minted}
\end{english}
\end{boxSolution}
\fi
\filbreak

% ==============================================================
% السؤال 8: حساب مدة حدث بالدقائق من وقتين بصيغة HHMM
% ==============================================================
\item
اكتب برنامجًا يستقبل وقتين (وقت بداية حدث ووقت نهايته) كأعداد صحيحة من $4$ منازل بنظام $24$ ساعة (مثلاً $1430$ تعني الساعة $14$:$30$)، ويحسب ويطبع مدة الحدث بالدقائق الإجمالية.
\ifdetailed
\begin{boxExample}
\begin{english}
\begin{minted}{text}
Please enter start time (HHMM):
1430
Please enter end time (HHMM):
1615
Duration in minutes = 105
\end{minted}
\end{english}
\end{boxExample}
\fi
\ifwithsols
\begin{boxSolution}
\begin{english}
\begin{minted}{csharp}
int start = int.Parse(Console.ReadLine());
int end = int.Parse(Console.ReadLine());

int startHours = start / 100;
int startMinutes = start % 100;
int startTotalMins = startHours * 60 + startMinutes;

int endHours = end / 100;
int endMinutes = end % 100;
int endTotalMins = endHours * 60 + endMinutes;

int duration = endTotalMins - startTotalMins;
Console.WriteLine("Duration in minutes = " + duration);
\end{minted}
\end{english}
\end{boxSolution}
\fi
\filbreak

\end{enumerate}

\vspace{1cm}
\begin{flushleft}
أرجو لكم وقتًا ممتعًا.

الأستاذ محمود اغبارية.
\end{flushleft}

\end{document}
